\documentclass[11pt,letterpaper]{article}

\usepackage[margin=0.6in]{geometry}
\usepackage{enumitem}
\usepackage{amsmath}

\pagestyle{empty}
\setlength{\parindent}{0pt}

\begin{document}

\begin{center}
    {\Large \textbf{16.07 Dynamics --- Quiz 5: Error Analysis and Quaternions}}
\end{center}

\vspace{4pt}

Name: \rule{3.2in}{0.4pt}

\vspace{14pt}
\begin{enumerate}[label=\textbf{\arabic*.}, leftmargin=*, itemsep=4pt, series=quiz]

\item To measure the period $P$ of the bifilar pendulum, you could either (i) time 10 consecutive oscillations in a single interval and divide by 10, or (ii) time a single oscillation 10 separate times and average the results. If each timed interval has an uncertainty of $\pm 0.1$~s, what is the uncertainty in $P$ for each approach? Which approach is better?
\vspace{2.0in}

\item A rotation matrix $\mathbf{Q}$ has 9 entries, and a quaternion $\mathbf{q}$ has 4. How many independent constraints must each satisfy to represent a valid rotation? How many degrees of freedom remain in each case?
\vspace{2.0in}

\item Because of numerical drift, the quaternion in your simulation has norm $\lVert\mathbf{q}\rVert = 1.01$. If you compute the rotation matrix $\mathbf{Q}(\mathbf{q}) = H^{T}L(\mathbf{q})\,R(\mathbf{q})^{T}H$ and apply it to a vector, what happens to that vector? How would you fix this? \emph{Hint:} $L(\mathbf{q})^{T}L(\mathbf{q}) = R(\mathbf{q})^{T}R(\mathbf{q}) = \lVert\mathbf{q}\rVert^2\,\mathbf{I}_4$.
\vspace*{\fill}

\end{enumerate}

\end{document}
